% Part: first-order-logic
% Chapter: sequent-calculus
% Section: proof-theoretic-notions

\documentclass[../../../include/open-logic-section]{subfiles}

\begin{document}

\begin{editorial}
  यह अनुभाग स्वाभाविक निगमन के सिद्धता-संबंध और संगतता की परिभाषाएँ एकत्र
  करता है।
\end{editorial}

\iftag{FOL}
      {\olfileid[hi]{fol}{seq}{ptn}}
      {\olfileid[hi]{pl}{seq}{ptn}}

\olsection{प्रमाण-सैद्धांतिक धारणाएँ}

\begin{explain}
जिस प्रकार हमने कई महत्त्वपूर्ण अर्थवैज्ञानिक धारणाएँ (वैधता, तार्किक
निष्पत्ति, संतुष्टनीयता) परिभाषित की हैं, उसी प्रकार अब उनकी संगत
\emph{प्रमाण-सैद्धांतिक धारणाएँ} परिभाषित करते हैं। इनकी परिभाषा में यह
सहारा नहीं लिया जाता कि !!{sentence}s के लिए कौन-सी !!{structure}s उन्हें
संतुष्ट करती हैं, बल्कि कुछ अनुक्रमों की !!{derivability} या
!!{nonderivability} का सहारा लिया जाता है। इन धारणाओं का एक-दूसरे से मेल
खाना एक महत्त्वपूर्ण खोज थी। यह मेल \emph{सुदृढ़ता} और \emph{पूर्णता प्रमेय}
का विषय है।
\end{explain}

\begin{defn}[प्रमेय]
!!^a{sentence}~$!A$ \emph{प्रमेय} है यदि~$\Log{LK}$ में अनुक्रम
$\quad \Sequent !A$ की कोई !!a{derivation} हो। तब $\Proves
!A$ लिखते हैं; यदि $!A$ प्रमेय नहीं है तो $\Proves/ !A$ लिखते हैं।
\end{defn}

\begin{defn}[!!^{derivability}]
!!^a{sentence}~$!A$ को किसी समुच्चय से \emph{!!{derivable}} कहते हैं; यदि वह
!!{sentence}s का समुच्चय~$\Gamma$ हो, तो इसे $\Gamma \Proves !A$ लिखते हैं।
यह तभी और केवल तभी होता है जब कोई परिमित उपसमुच्चय
~$\Gamma_0 \subseteq \Gamma$ और कोई क्रम $\Gamma_0'$ ऐसा हो जो~$\Gamma_0$
के !!{sentence}s से बना हो, तथा $\Log{LK}$
$\Gamma_0' \Sequent !A$ को !!{derive} करे। यदि $!A$, $\Gamma$ से
!!{derivable} नहीं है तो $\Gamma \Proves/ !A$ लिखते हैं।
\end{defn}

संकुचन, दुर्बलीकरण और विनिमय नियमों के कारण~$\Gamma_0'$ में !!{sentence}s
का क्रम और संख्या महत्त्व नहीं रखते: यदि अनुक्रम $\Gamma_0' \Sequent !A$
!!{derivable} है, तो $\Gamma_0'' \Sequent !A$ भी हर ऐसे $\Gamma_0''$ के लिए
वैसा ही है जिसका निर्माण~$\Gamma_0'$ वाले !!{sentence}s से हुआ हो। उदाहरण के लिए, यदि
$\Gamma_0 = \{!B, !C\}$ हो, तो $\Gamma_0' = \tuple{!B, !B, !C}$ और $\Gamma_0'' =
\tuple{!C, !C, !B}$, दोनों ऐसे क्रम हैं जो~$\Gamma_0$ के ही
!!{sentence}s से बने हैं। यदि इनमें से एक क्रम वाला अनुक्रम !!{derivable} है, तो
दूसरा भी है; जैसे:
\begin{prooftree}
  \AxiomC{}
  \Deduce$!B, !B, !C \fCenter !A$
  \RightLabel{\LeftR{\Contraction}}
  \UnaryInf$!B, !C \fCenter !A$
  \RightLabel{\LeftR{\Exchange}}
  \UnaryInf$!C, !B \fCenter !A$
  \RightLabel{\LeftR{\Weakening}}
  \UnaryInf$!C, !C, !B \fCenter !A$
\end{prooftree}
अब से, यदि $\Gamma_0$, !!{sentence}s का परिमित समुच्चय है, तो
$\Gamma_0 \Sequent !A$ से हमारा आशय ऐसे किसी भी अनुक्रम से होगा जिसका
पूर्वपक्ष~$\Gamma_0$ के !!{sentence}s का क्रम हो; आवश्यकता पड़ने पर संकुचन,
विनिमय और दुर्बलीकरण को मौन रूप से शामिल मानेंगे।

\begin{defn}[संगतता]
वाक्यों का समुच्चय~$\Gamma$ \emph{असंगत} है तभी और केवल तभी जब कोई परिमित
उपसमुच्चय~$\Gamma_0 \subseteq \Gamma$ ऐसा हो कि $\Log{LK}$
$\Gamma_0 \Sequent \quad$ को !!{derive} करे। यदि $\Gamma$ असंगत नहीं है,
अर्थात् हर परिमित $\Gamma_0 \subseteq \Gamma$ के लिए $\Log{LK}$
$\Gamma_0 \Sequent \quad$ को !!{derive} नहीं करता, तो उसे \emph{संगत} कहते
हैं।
\end{defn}

\begin{prop}[स्वतुल्यता]
\ollabel{prop:reflexivity}
यदि $!A \in \Gamma$, तो $\Gamma \Proves !A$।
\end{prop}

\begin{proof}
प्रारंभिक अनुक्रम $!A \Sequent !A$ !!{derivable} है और $\{!A\}
\subseteq \Gamma$।
\end{proof}
  
\begin{prop}[एकदिष्टता]
\ollabel{prop:monotonicity}
यदि $\Gamma \subseteq \Delta$ और $\Gamma \Proves !A$, तो $\Delta
\Proves !A$।
\end{prop}

\begin{proof}
मान लें $\Gamma \Proves !A$; अर्थात् कोई परिमित $\Gamma_0
\subseteq \Gamma$ ऐसा है कि $\Gamma_0 \Sequent !A$ !!{derivable} है। चूँकि
$\Gamma \subseteq \Delta$, इसलिए $\Gamma_0$,~$\Delta$ का भी परिमित
उपसमुच्चय है। अतः $\Gamma_0
\Sequent !A$ की !!{derivation} यह भी दिखाती है
कि $\Delta \Proves !A$।
\end{proof}

\begin{prop}[संक्रामकता]
\ollabel{prop:transitivity}
यदि $\Gamma \Proves !A$ और $\{!A\} \cup \Delta \Proves
!B$, तो $\Gamma \cup \Delta \Proves !B$।
\end{prop}

\begin{proof}
यदि $\Gamma \Proves !A$, तो कोई परिमित $\Gamma_0 \subseteq \Gamma$ और कोई
!!a{derivation}~$\pi_0$ ऐसी है जिसका अंतिम अनुक्रम $\Gamma_0 \Sequent !A$ है। यदि $\{!A\}
\cup \Delta \Proves !B$, तो किसी परिमित उपसमुच्चय $\Delta_0
\subseteq \Delta$ के लिए कोई !!a{derivation}~$\pi_1$ ऐसी है जिसका अंतिम अनुक्रम $!A, \Delta_0
\Sequent !B$ है। निम्नलिखित !!{derivation} पर विचार करें:
\begin{prooftree}
  \AxiomC{}
  \RightLabel{$\pi_0$}
  \Deduce$\Gamma_0 \fCenter !A$
  \AxiomC{}
  \RightLabel{$\pi_1$}
  \Deduce$!A, \Delta_0 \fCenter !B$
  \RightLabel{\Cut}
  \BinaryInf$\Gamma_0, \Delta_0 \fCenter !B$
\end{prooftree}
चूँकि $\Gamma_0 \cup \Delta_0 \subseteq \Gamma \cup \Delta$, इससे
$\Gamma \cup \Delta \Proves !B$ सिद्ध होता है।
\end{proof}

ध्यान दें कि विशेषतः इसका अर्थ है: यदि $\Gamma \Proves !A$ और $!A
\Proves !B$, तो $\Gamma \Proves !B$। इससे यह भी मिलता है कि यदि $!A_1,
\dots, !A_n \Proves !B$ और $\Gamma \Proves !A_i$ प्रत्येक~$i$ के लिए, तो
$\Gamma \Proves !B$।

\begin{prop}
\ollabel{prop:incons}
$\Gamma$ असंगत है तभी और केवल तभी जब $\Gamma \Proves {!A}$ हर
  वाक्य~$!A$ के लिए सत्य हो।
\end{prop}

\begin{proof}
अभ्यास।
\end{proof}

\begin{prob}
\olref[fol][seq][ptn]{prop:incons} सिद्ध कीजिए।
\end{prob}

\begin{prop}[संहतता]
  \ollabel{prop:proves-compact}
  \begin{enumerate}
  \item यदि $\Gamma \Proves !A$, तो कोई परिमित उपसमुच्चय $\Gamma_0
    \subseteq \Gamma$ ऐसा है कि $\Gamma_0 \Proves !A$।
  \item यदि~$\Gamma$ का हर परिमित उपसमुच्चय संगत है, तो $\Gamma$ संगत है।
  \end{enumerate}
\end{prop}

\begin{proof}
  \begin{enumerate}
    \item यदि $\Gamma \Proves !A$, तो कोई परिमित उपसमुच्चय
      $\Gamma_0 \subseteq \Gamma$ ऐसा है कि अनुक्रम $\Gamma_0
      \Sequent !A$ की कोई !!a{derivation} है। फलतः $\Gamma_0
      \Proves !A$।
    \item यदि $\Gamma$ असंगत है, तो कोई परिमित उपसमुच्चय
      ~$\Gamma_0 \subseteq \Gamma$ ऐसा है कि $\Log{LK}$, $\Gamma_0 \Sequent \quad$
      को !!{derive} करता है। किंतु तब $\Gamma_0$,~$\Gamma$ का असंगत परिमित
      उपसमुच्चय है।
  \end{enumerate}
\end{proof}

\end{document}
